\section{Work a checksum by hand}
The checksum rule requires the sum of the header, payload, and trailer to be a
multiple of 256. Given all preceding bytes with sum \(S\), choose trailer
\[
 c=(-S)\bmod256,\qquad 0\leq c\leq255.
\]
This is an additive checksum, not encryption or authentication. It detects some
accidental changes. It does not prove that the sender is trustworthy or that each
field has been decoded correctly.

\begin{MARSWorkedExample}{The default twelve-byte packet}
Read all numbers in this byte sequence as hexadecimal:
\begin{quote}\texttt{4d 50 01 01 00 03 12 34 7f 80 ff 1a}\end{quote}
The first eight bytes sum to 232 in decimal. The three payload bytes sum to
\(127+128+255=510\), whose residue modulo 256 is 254. Header and payload together
sum to 742, with residue 230. The trailer \texttt{1a} is 26; the complete sum is
768, exactly three times 256. The returned consumed length is 12. The output words
are flags 1, sequence 4660, payload length 3, and payload-sum residue 254.
\end{MARSWorkedExample}
The header and payload need separate interpretations. The checksum uses the raw
byte values. Adding decoded sequence number 4660 in place of its two bytes 18 and
52 changes the calculation. The record's payload sum excludes both header and
trailer. A single accumulator for the entire packet does not directly provide that
separate output field.

At most 73 packet bytes contribute to the full checksum, giving a loose unsigned
sum bound of \(73\cdot255=18615\). The payload-only sum is at most 16320. Both fit
comfortably in a signed 32-bit word, so the example can defer masking until the
end. Applying \texttt{andi} with 255 retains the low eight bits and computes the
nonnegative residue for these sums.

Signed payload-byte loads provide a useful counterexample to simplistic testing.
For a high-bit byte, \texttt{lb} differs from \texttt{lbu} by exactly 256. Their sums
therefore have the same residue modulo 256. A parser using signed payload loads
and masking the final sums can satisfy this contract. A signed low byte used in
the sequence-field OR generally does not: its extended high bits corrupt the
decoded integer. Check field values and checksum values independently.
